% !TeX root = ../tesis.tex

% =============================================================================
%  SECCIÓN 3.1.1 — Definición de Conceptos Fundamentales: Introducción y
%                   Glosario de Términos Instrumentales
% =============================================================================

\section{Definición de Conceptos Fundamentales}
\label{sec:conceptos-fundamentales}

El análisis de redes de transporte público desde un enfoque computacional exige
fijar con precisión los términos que articulan el diagnóstico del sistema y el
diseño de intervenciones. Este capítulo cumple dos propósitos complementarios:
primero, establecer el glosario de conceptos instrumentales que la investigación
emplea a lo largo de todos sus capítulos ---desde los estándares de datos
abiertos hasta la representación matemática de la red---; segundo, desarrollar
los cuatro conceptos teórico-analíticos que sustentan la metodología
---vulnerabilidad urbana, planeación estratégica, eficiencia urbana e
interconectividad entre redes--- con base en literatura especializada en teoría
de grafos, redes complejas y análisis de movilidad urbana. Ningún término
técnico que aparezca en los capítulos posteriores llega al lector sin una
definición formal previa.

% -----------------------------------------------------------------------
\subsection{Glosario de términos instrumentales}
\label{subsec:glosario-instrumental}
% -----------------------------------------------------------------------

Los siguientes términos constituyen el vocabulario común de toda la
investigación. Se presentan antes de los desarrollos teóricos de las secciones
subsecuentes porque son condición de lectura de los indicadores, los algoritmos
y los resultados.

\subsubsection{Sistema de Información Geográfica (\gis{})}
\label{subsubsec:glosario-gis}
Un \gis{} es un sistema informático que captura, almacena, verifica y despliega
datos vinculados a posiciones en la superficie terrestre
\citep{longley2015}. Su aplicación al análisis de redes de transporte permite
integrar capas de información espacial heterogéneas ---geometría de paradas,
trazados de rutas, población censal y red vial--- para calcular indicadores de
cobertura y accesibilidad sobre un mismo marco georreferenciado
\citep{ugur2025}. En esta investigación, el \gis{} se implementa con
herramientas de código abierto ---PostGIS, GeoPandas y QGIS---, cuya cadena
de procesamiento produce los indicadores definidos en
\autoref{sec:indicadores-formulas}.

\subsubsection{Información Geográfica Voluntaria (\textsc{vgi})}
\label{subsubsec:glosario-vgi}
El término Información Geográfica Voluntaria ---\textsc{vgi}, del inglés
\textit{Volunteered Geographic Information}--- refiere a los datos geoespaciales
producidos y publicados colectivamente por ciudadanos en plataformas abiertas
en línea \citep{haklay2010good}. En esta investigación, la red vial de base
procede de fuentes oficiales ---el portal de datos abiertos de la \cdmx{} y
el \textsc{inegi}---; OpenStreetMap (\textsc{osm}) opera como fuente secundaria
para cubrir los vacíos donde la cobertura oficial es incompleta, función que
resulta indispensable en la construcción de los grafos de acceso peatonal
sobre los que se calculan las isócronas del Indicador de Cobertura $C$
(§\ref{subsec:indicador-cobertura}).

Su completitud supera el 80\,\% a escala global, con valores especialmente
altos en ciudades latinoamericanas de la talla de la \zmvm{}; la precisión
posicional en zonas urbanas densas no difiere significativamente de la de
fuentes cartográficas institucionales \citep{barrington2017world,
haklay2010good}, lo que sustenta su uso como complemento de la cartografía
oficial en este análisis.

\subsubsection{Estándar de datos abiertos de transporte (\gtfs{})}
\label{subsubsec:glosario-gtfs}
El \gtfs{} (\textit{General Transit Feed Specification}) es el formato de
archivos de texto plano mediante el cual los operadores de transporte público
publican rutas, paradas, frecuencias y horarios de forma interoperable
\citep{fortin2016innovative}. Creado en 2005 por Google y la agencia TriMet
de Portland, Oregón, hoy lo adoptan más de 10\,000 operadoras en todo el
mundo. En esta investigación, los feeds \gtfs{} de todos los modos de la
red de movilidad integrada de la \cdmx{} ---disponibles en el portal de
datos abiertos de la ciudad \citep{adip_gtfs_2024}--- constituyen el insumo
primario que el Apímetro descarga e integra para construir el grafo
multimodal de análisis (§\ref{sec:apimetro-intro}).

\subsubsection{Normalización de datos de transporte}
\label{subsubsec:glosario-normalizacion}
La normalización de datos es el proceso mediante el cual registros de
distintas fuentes se reorganizan en un esquema relacional canónico, libre de
redundancias y anomalías de actualización \citep{ramakrishnan2003}. En el
contexto de esta investigación, cada sistema de transporte de la \zmvm{}
opera bajo un orden de gobierno diferente ---federal, estatal o
municipal--- y genera sus propios feeds \gtfs{} con convenciones de
nomenclatura, codificación y granularidad temporal distintas. El Apímetro
aplica este proceso sobre \textsc{postgresql} \citep{postgresql2024}, lo
que transforma esos feeds heterogéneos en un esquema unificado; condición
indispensable para construir el grafo multimodal $\grafo$ descrito en el
\autoref{ch:Capas-codigo}.

\subsubsection{Grafo de transporte ponderado y dirigido}
\label{subsubsec:glosario-grafo}
Un grafo de transporte es la representación matemática de la red en la que
cada parada o estación corresponde a un nodo $v \in \nodos$ y cada conexión
operada entre dos paradas es una arista dirigida $e \in \aristas$
\citep{newman2010networks}. La función de peso $W: \aristas \to \mathbb{R}^+$
asigna a cada arista el costo compuesto de recorrerla ---tiempo de viaje,
tiempo de espera y, cuando aplica, penalización por transbordo---, de modo
que la estructura $\grafo = \{\nodos,\, \aristas,\, W\}$ sirve de base para
calcular los indicadores del \autoref{sec:indicadores-formulas} mediante el
algoritmo de \dijkstra{}. La formalización matemática completa, junto con sus
propiedades topológicas e implementación computacional, se desarrolla en la
\autoref{sec:grafos-matematicos}.

\subsubsection{Red de transporte multimodal de la \zmvm{}}
\label{subsubsec:glosario-red-multimodal}
La red de transporte multimodal de la \zmvm{} comprende doce sistemas de
transporte estructurado ---Metro, Metrobús, \brt{}, Tren Suburbano, Cablebús,
Mexicable, Mexibús, Trolebús, Tren Ligero, Tren Interurbano, \textsc{rtp} y
rutas de superficie concesionadas--- integrados bajo una estructura unificada
de análisis \citep{semovi2020diagnostico}. En términos de grafo, las aristas
de trayecto representan los recorridos dentro de cada línea; las aristas de
transferencia conectan modos distintos en el mismo punto físico, con un peso
que refleja el tiempo de caminata entre andenes y el tiempo de espera según la
frecuencia del servicio \citep{molinero2002}.

\subsubsection{\textit{Hub} y Centro de Transferencia Modal (\textsc{cetram})}
\label{subsubsec:glosario-hub-cetram}
En la teoría de redes, un \termino{hub} es un nodo con un grado de conexión
excepcionalmente alto respecto a la distribución del resto del sistema
\citep{newman2010networks}. En el transporte metropolitano de la \cdmx{},
los Centros de Transferencia Modal (\textsc{cetram}) ---Pantitlán, Indios
Verdes, Cuatro Caminos y Taxqueña--- operan como \textit{hubs} que
concentran la articulación entre las redes capilares de superficie y las
arterias masivas \citep{semovi2020diagnostico}. Su alta Centralidad de
Intermediación $B(v)$ los convierte, simultáneamente, en los puntos de mayor
vulnerabilidad ante fallas o saturación del sistema
\citep{lotero2014vulnerabilidad} (§\ref{subsec:centralidad-intermediacion}).

\subsubsection{Corredor de alta capacidad anillar}
\label{subsubsec:glosario-corredor-anillar}
Bajo la categoría de corredor de alta capacidad se agrupan líneas de
transporte público con infraestructura exclusiva ---carril confinado,
plataformas de abordaje a nivel, control centralizado de frecuencias---
que alcanzan rendimientos operativos comparables a los sistemas ferroviarios
a una fracción de su costo \citep{molinero2002}. La variante anillar se distingue de los trazados radiales
por su geometría: en lugar de converger al centro, el corredor describe un arco
concéntrico paralelo a la frontera entre el área central y la primera corona
periférica, articulando cuadrantes de la periferia entre sí y desviando
los flujos que de otro modo saturarían los nodos de alta intermediación
\citep{aldous2019optimal}. Dos proyectos recientes ilustran la viabilidad
del modelo: el orbital ferroviario de 90\,km en Melbourne
\citep{suburbanrailloop2021} y la Gran Línea Circular de 54 estaciones
inaugurada en Moscú en 2023 \citep{efe2023}. Para la \zmvm{}, cuya periferia
carece de vínculos directos de transporte estructurado entre cuadrantes
\citep{obras2025}, esta investigación evalúa el impacto hipotético de
implementar corredores anillares en los cuatro cuadrantes del Anillo
Periférico de la \cdmx{} como escenario de demostración metodológica.

\subsubsection{Términos de arquitectura de software}
\label{subsubsec:glosario-software}

Los siguientes términos corresponden a estándares y patrones de ingeniería
de software empleados en la implementación del Apímetro. Se definen aquí
para que el lector no especializado pueda seguir la descripción técnica del
\autoref{ch:Capas-codigo} sin ambigüedad.

\paragraph{JSON (\textit{JavaScript Object Notation}).}
Formato ligero de intercambio de datos basado en texto plano,
independiente del lenguaje de programación y legible por humanos
\citep{bray2017json}. Es la representación estándar para la comunicación
entre servicios web modernos.

\paragraph{GeoJSON.}
Extensión del formato JSON para la codificación de objetos geográficos
---puntos, líneas y polígonos--- con sus atributos espaciales y temáticos
\citep{butler2016geojson}. El Apímetro entrega todos sus resultados en
este formato.

\paragraph{\api{} \textsc{rest}.}
Estilo arquitectónico para sistemas de software distribuidos que define
seis restricciones ---entre ellas la interfaz uniforme y la ausencia de
estado--- que, aplicadas a servicios web sobre \textsc{http}, producen
interfaces predecibles e interoperables \citep{fielding2000rest}.

\paragraph{Arquitectura hexagonal (\textit{Ports and Adapters}).}
Patrón de diseño que aísla la lógica central de la aplicación
---el \textit{dominio}--- de cualquier tecnología de entrada/salida
mediante puertos abstractos e implementaciones intercambiables
\citep{cockburn2005hexagonal}. Adoptado en el Apímetro para separar
las reglas de procesamiento de grafos de sus adaptadores de
persistencia y exposición de datos.

\paragraph{Go (lenguaje de programación).}
Lenguaje compilado y tipado estáticamente desarrollado en Google,
diseñado para la concurrencia estructurada y la compilación rápida
\citep{go2024spec}. El Apímetro está implementado en Go por su
rendimiento en operaciones de entrada/salida intensiva y su soporte
nativo para concurrencia mediante \textit{goroutines}.

\paragraph{Pipeline (cadena metodológica).}
Secuencia ordenada de etapas de procesamiento en la que la salida de
cada fase es la entrada de la siguiente \citep{ramakrishnan2003}. En
este trabajo, el término designa la cadena Apímetro $\to$ VFTModel $\to$
motor de indicadores que transforma feeds \gtfs{} crudos en resultados
de análisis urbano.

\subsubsection{Herramientas de software de la investigación}
\label{subsubsec:glosario-herramientas}

Esta investigación desarrolló tres herramientas de código abierto que
articulan la cadena metodológica completa. Se presentan aquí en una línea
cada una; su arquitectura técnica se documenta en el
\autoref{ch:Capas-codigo}.

\paragraph{Apímetro.}
Motor de adquisición e integración de datos: descarga, normaliza y unifica
feeds \gtfs{} de múltiples operadores en el esquema relacional de
\textsc{postgresql}, produciendo el archivo Geojson multimodal $\grafo$ que alimenta
el análisis (§\ref{sec:apimetro-intro}).

\paragraph{VFTModel.}
Modelo de análisis urbano: recibe el grafo $\grafo$ y calcula los
indicadores de cobertura, eficiencia, centralidad y robustez definidos en
la \autoref{sec:indicadores-formulas}; implementado con arquitectura
hexagonal para separar la lógica analítica de sus adaptadores de datos.

\paragraph{Transport-\gis{}.}
Capa de visualización geoespacial: transforma los resultados del VFTModel
en capas cartográficas comparables entre escenarios base e hipotéticos,
exportadas en formato GeoJSON para su publicación en el portal de datos
abiertos.

\medskip

Los términos definidos en este glosario aparecen de forma recurrente en los
capítulos metodológicos y de resultados. Que figuren aquí antes de cualquier
desarrollo analítico garantiza que los Objetivos Específicos del 1 al 5
---que los invocan directamente en su formulación
(§\ref{subsec:obj-especificos})--- puedan leerse sin ambigüedad desde el
inicio del documento.
